\startmode[single]

% \enabletrackers[publications, publications.crossref, publications.details, publications.cite, publications.strings]
%
% \usebtxdataset[main][bib/default.bib,bib/physik.bib,bib/mathe.bib,bib/seminarfach.bib]
% \usebtxdefinitions[aps]
% \setupbtx[dataset=main]
% \definebtxrendering[bibrendering][aps][dataset=main]

\setvariables[meta]
  [title={Elektronenbeugunsröhre},
   subtitle={Versuchsprotokoll},
   date={2024-10-18},
   author={Niklas von Hirschfeld},
  ]

\setupinteraction[title={\getvariable{meta}{title}}]

\starttext

\startfrontmatter
\component c_titlepage
\component c_toc
\stopfrontmatter

\setuppagenumber[number=1]

\startbodymatter

\startchapter[title={Kernphysik 04 Abschirmung und Ausbreitung}]

\stopmode


\section{A1}

\startplacetable[location={here},title={Ablesen der Zeiträume}]
\setupTABLE[c][1][background=color,backgroundcolor=gray]
\setupTABLE[r][1][background=color,backgroundcolor=gray]
\bTABLE[frame=off,offset=1mm]
\startCSV
Zeitabschnitt, 1, 2, 3, 4, Durchschnittdifferenz
Halbwertsdicke in $mm$, 6, 12, 17, 24, 6
\stopCSV
\eTABLE
\stopplacetable

\startframedtipp
Beim grafischen Ablesen auch {\bf MARKIEREN}
\stopframedtipp

\startframedtipp
+ANKI
\startformula
a^{-b}=\frac{1}{a^b}
\stopformula

\blank
\hrule
\blank

\startformula
\ln(a)=-\ln(\frac{1}{a})
\stopformula
\stopframedtipp




\startmode[single]

\stopchapter

\stopbodymatter

\startbackmatter
% \component c_definitions
\component c_bibliography
\stopbackmatter
\stoptext

\stopmode
